\section{Ekspansi State dalam Basis Eigen: Mekanisme Konvergensi Kuantum}\label{appendix-g}


\hypertarget{postulat-dasar-spektrum-operator-hermitian}{%
\subsection{Postulat Dasar: Spektrum Operator Hermitian}\label{postulat-dasar-spektrum-operator-hermitian}}

Dalam mekanika kuantum, Hamiltonian \(\hat{H}\) adalah operator linear Hermitian. Berdasarkan \textbf{Teorema Spektral} dalam aljabar linear:

\textbf{Aksioma \& Konvensi:} 1. \textbf{Eigensystem:} Terdapat sekumpulan nilai eigen riil \(\{E_n\}\) dan vektor eigen \(\{|v_n\rangle\}\) yang memenuhi persamaan: \[ \hat{H} |v_n\rangle = E_n |v_n\rangle \] 2. \textbf{Ortonormalitas:} Vektor-vektor eigen ini bersifat ortonormal, yaitu \(\langle v_m | v_n \rangle = \delta_{mn}\). 3. \textbf{Kelengkapan (Completeness):} Vektor-vektor eigen \(\{|v_n\rangle\}\) membentuk basis lengkap untuk ruang Hilbert \(\mathcal{H}\).

\textbf{Filosofi:} Basis eigen adalah ``bahasa alami'' dari Hamiltonian. Sama seperti warna putih yang dapat diuraikan menjadi spektrum warna pelangi, sebuah state kuantum abstrak \(|\psi\rangle\) dapat diuraikan menjadi komponen-komponen energi yang menyusunnya.

\begin{center}\rule{0.5\linewidth}{0.5pt}\end{center}

\hypertarget{jembatan-1-prinsip-superposisi-dan-resolusi-identitas}{%
\subsection{Jembatan 1: Prinsip Superposisi dan Resolusi Identitas}\label{jembatan-1-prinsip-superposisi-dan-resolusi-identitas}}

Karena \(\{|v_n\rangle\}\) adalah basis lengkap, kita dapat menyatakan sembarang state \(|\psi\rangle\) sebagai kombinasi linear (superposisi) dari vektor-vektor eigen tersebut.

\textbf{Alat (Resolusi Identitas):} Sebuah identitas operator \(\hat{I}\) dapat dinyatakan sebagai jumlah dari semua operator proyeksi ke basis eigen: \[ \hat{I} = \sum_n |v_n\rangle\langle v_n | \]

\textbf{Jembatan Gagasan:} Dengan menyisipkan identitas ke dalam \(|\psi\rangle\), kita memproyeksikan state tersebut ke seluruh spektrum Hamiltonian: \[ |\psi\rangle = \hat{I} |\psi\rangle = \sum_n |v_n\rangle \langle v_n | \psi \rangle \] \[ |\psi\rangle = \sum_n \alpha_n |v_n\rangle \] di mana \(\alpha_n = \langle v_n | \psi \rangle\) adalah \textbf{amplitudo probabilitas} sistem berada pada tingkat energi \(E_n\).

\begin{center}\rule{0.5\linewidth}{0.5pt}\end{center}

\hypertarget{jembatan-2-normalisasi-dan-konservasi-probabilitas}{%
\subsection{Jembatan 2: Normalisasi dan Konservasi Probabilitas}\label{jembatan-2-normalisasi-dan-konservasi-probabilitas}}

Agar \(|\psi\rangle\) merepresentasikan sistem fisik yang nyata, total probabilitas harus bernilai satu.

\textbf{Derivasi Runtut:} \[ \langle \psi | \psi \rangle = \left( \sum_m \alpha_m^* \langle v_m | \right) \left( \sum_n \alpha_n |v_n\rangle \right) \] Karena sifat ortonormalitas (\(\langle v_m | v_n \rangle = \delta_{mn}\)), maka: \[ \sum_n |\alpha_n|^2 = 1 \]

\textbf{Jembatan Gagasan:} Koefisien \(|\alpha_n|^2\) (berdasarkan Aturan Born) menunjukkan bobot kontribusi setiap tingkat energi terhadap state total. Jika \(|\alpha_0|^2\) besar, maka state tersebut sangat mirip dengan \emph{ground state}.

\begin{center}\rule{0.5\linewidth}{0.5pt}\end{center}

\hypertarget{jembatan-3-nilai-ekspektasi-dalam-representasi-spektral}{%
\subsection{Jembatan 3: Nilai Ekspektasi dalam Representasi Spektral}\label{jembatan-3-nilai-ekspektasi-dalam-representasi-spektral}}

Bagaimana ekspansi ini membantu kita menghitung energi rata-rata?

\textbf{Derivasi Energi:} \[ \langle E \rangle = \langle \psi | \hat{H} | \psi \rangle = \langle \psi | \hat{H} \left( \sum_n \alpha_n |v_n\rangle \right) \] Karena \(\hat{H} |v_n\rangle = E_n |v_n\rangle\): \[ \langle E \rangle = \left( \sum_m \alpha_m^* \langle v_m | \right) \left( \sum_n \alpha_n E_n |v_n\rangle \right) \] \[ \langle E \rangle = \sum_n E_n |\alpha_n|^2 \]

\textbf{Filosofi Jembatan:} Energi yang kita ukur (nilai ekspektasi) adalah \textbf{rata-rata tertimbang} dari seluruh spektrum energi Hamiltonian, di mana bobotnya adalah probabilitas \(|\alpha_n|^2\) dari ekspansi state kita.

\begin{center}\rule{0.5\linewidth}{0.5pt}\end{center}

\hypertarget{aplikasi-pada-vqe-mekanisme-konvergensi}{%
\subsection{Aplikasi pada VQE: Mekanisme Konvergensi}\label{aplikasi-pada-vqe-mekanisme-konvergensi}}

Dalam VQE, kita memiliki state parametrik \(|\psi(\theta)\rangle\). Berarti amplitudo probabilitasnya juga bergantung pada parameter: \(\alpha_n(\theta)\).

\textbf{Logika Optimasi:} 1. \textbf{Keadaan Awal:} Saat \(\theta\) acak, \(|\psi(\theta)\rangle\) memiliki komponen di banyak tingkat energi (banyak \(\alpha_n\) yang tidak nol). Energi \(\langle E \rangle\) akan bernilai tinggi. 2. \textbf{Proses Minimisasi:} Algoritma optimasi mengubah \(\theta\) untuk menurunkan \(\langle E \rangle\). 3. \textbf{Filter Energi:} Secara matematis, proses ini ``membuang'' komponen state eksitasi (\(n > 0\)) dengan cara mengecilkan \(|\alpha_n(\theta)|^2\) dan memperbesar \(|\alpha_0(\theta)|^2\).

\textbf{Limit Konvergensi:} \[ \lim_{\theta \to \theta^*} \alpha_n(\theta) = \begin{cases} 1 & \text{untuk } n=0 \\ 0 & \text{untuk } n \ne 0 \end{cases} \] \[ \lim_{\theta \to \theta^*} \langle E \rangle = E_0 \]

\begin{center}\rule{0.5\linewidth}{0.5pt}\end{center}

\hypertarget{kesimpulan}{%
\subsection{Kesimpulan}\label{kesimpulan-g}}

Ekspansi state adalah alat yang memungkinkan kita melihat ``jeroan'' dari sirkuit kuantum. Ia membuktikan bahwa menurunkan nilai ekspektasi energi secara numerik setara dengan \textbf{menyelaraskan sirkuit kuantum} kita agar tepat berada pada basis eigen energi terendah (\emph{ground state}). Tanpa konsep ekspansi ini, kita tidak akan memiliki justifikasi matematis mengapa VQE dapat menemukan solusi portfolio optimal.
